\section{Постановка задачи}

\subsection{Языковая последовательность}

Пусть \((\Omega,\F,\Prob)\) --- вероятностное пространство. Речь задается последовательностью
\begin{equation}
  X_1,X_2,\ldots,X_T, \qquad X_t\in\A,
  \label{eq:language_sequence}
\end{equation}
где \(\A=\{1,\ldots,V\}\) --- конечный словарь языковых единиц. Наблюденный префикс до позиции \(t\) задает \(\sigma\)-алгебру
\begin{equation}
  \F^X_t=\sigma(X_1,\ldots,X_t).
\end{equation}
Истинный условный закон следующей единицы имеет вид
\begin{equation}
  p_t(a)=\Prob(X_{t+1}=a\mid \F^X_t), \qquad a\in\A.
  \label{eq:true_language_law}
\end{equation}
Этот закон неизвестен и оценивается моделью.

В рекурсивной схеме для каждой позиции \(t\) строится не одно, а несколько приближений к \(p_t\):
\begin{equation}
  Q_{t,0},Q_{t,1},\ldots,Q_{t,K}, \qquad Q_{t,k}\in\Delta^{V-1}.
  \label{eq:recursive_predictions}
\end{equation}
Здесь \(k=0\) обозначает начальное распределение, а \(k=1,\ldots,K\) --- результаты внутренних уточнений. Для фиксированной позиции можно остановиться на любом шаге от \(0\) до \(K\).

\subsection{Внутреннее состояние и доступная информация}

Рекурсивная модель хранит скрытое состояние \(Z_{t,k}\). Начальное состояние строится по префиксу:
\begin{equation}
  Z_{t,0}=E_{\phi}(X_1,\ldots,X_t),
  \label{eq:initial_state}
\end{equation}
где \(E_\phi\) --- исходный языковой преобразователь. Далее применяется оператор уточнения
\begin{equation}
  Z_{t,k+1}=U_\theta(Z_{t,k}, C_t, Q_{t,k}), \qquad k=0,\ldots,K-1.
  \label{eq:recursive_update}
\end{equation}
Величина \(C_t\) обозначает представление контекста, которое не меняется внутри рекурсивного цикла или меняется медленнее, чем \(Z_{t,k}\). После каждого шага скрытое состояние переводится в распределение на словаре:
\begin{equation}
  Q_{t,k}(a)=\frac{\exp s_{t,k}(a)}{\sum_{b\in\A}\exp s_{t,k}(b)}, \qquad a\in\A.
  \label{eq:softmax_distribution}
\end{equation}

Информация, доступная к внутреннему шагу \(k\), задается как
\begin{equation}
  \Hh_{t,k}=\F^X_t\vee\sigma(Z_{t,0},Q_{t,0},\ldots,Z_{t,k},Q_{t,k},S_{t,0},\ldots,S_{t,k}),
  \label{eq:information_flow}
\end{equation}
где \(S_{t,k}\) --- вектор вычисляемых статистик. В него можно включить энтропию распределения, отрыв первой вероятности от второй, изменение распределения между соседними шагами и норму изменения скрытого состояния. Эти статистики не вводят новую сущность задачи: они являются функциями уже построенных \(Z_{t,k}\) и \(Q_{t,k}\).

\subsection{Потери и оракульное время остановки}

После того как известна истинная следующая языковая единица \(X_{t+1}\), потери при остановке на шаге \(k\) определяются как
\begin{equation}
  L_{t,k}=-\log Q_{t,k}(X_{t+1}).
  \label{eq:nll_loss}
\end{equation}
Если учитывается не только качество прогноза, но и стоимость вычислений, вводится функция стоимости \(c(k)\). В простейшем случае \(c(k)=k\). Оракульное время остановки равно
\begin{equation}
  \tau_t^\star=\argmin_{0\le j\le K}\{L_{t,j}+\lambda c(j)\},
  \label{eq:oracle_stopping_time}
\end{equation}
где \(\lambda\ge0\) --- коэффициент цены вычислений. При равенстве нескольких минимумов выбирается наименьший \(j\), поскольку более ранняя остановка имеет ту же ошибку при меньшей или равной стоимости.

Оракульное время не является реализуемым правилом вывода: для его вычисления требуется знать все будущие внутренние шаги и истинный ответ. Оно используется как обучающая целевая величина и как статистический объект анализа.

\subsection{Основная задача предсказания распределения}

На шаге \(k\) модель времени остановки должна выдать условное распределение
\begin{equation}
  \widehat p_{t,k}(j)=\Prob_\theta(\tau_t^\star=j\mid\Hh_{t,k}), \qquad j=0,\ldots,K,
  \label{eq:main_distribution}
\end{equation}
Для пакета размера \(B\), длины последовательности \(T\) и \(K+1\) возможных шагов выход имеет размер
\begin{equation}
  B\times T\times (K+1)\times (K+1).
  \label{eq:language_output_size}
\end{equation}
Третья ось соответствует текущему шагу \(k\), четвертая --- возможному значению оракульного времени \(j\).

В~\eqref{eq:main_distribution} символ \(\Prob_\theta\) обозначает вероятностную меру, задаваемую параметрической моделью с параметром \(\theta\) (обучаемым вектором весов): \(\widehat p_{t,k}(j)\) --- \(\Hh_{t,k}\)-измеримая вероятность на \(\{0,\ldots,K\}\). Истинное условное распределение \(\tau_t^\star\) по \(\Hh_{t,k}\) относительно исходной меры \(\Prob\) в общем случае неизвестно и не совпадает с \(\Prob_\theta\). Событие \(\{\tau_t^\star=j\}\subset\Omega\) задается после наблюдения \(X_{t+1}\) и полной траектории внутренних прогнозов, достаточной для вычисления~\eqref{eq:oracle_stopping_time}.

Эта постановка принципиально отличается от прямого прогнозирования только индикатора \(\{\tau_t^\star=k\}\) без учета масс на \(\{\tau_t^\star<k\}\) и \(\{\tau_t^\star>k\}\). Модель сохраняет информацию о трех ситуациях:
\begin{align}
  P_{t,k}^{\mathrm{past}} &= \sum_{j<k}\widehat p_{t,k}(j), \\
  P_{t,k}^{\mathrm{now}} &= \widehat p_{t,k}(k), \\
  P_{t,k}^{\mathrm{future}} &= \sum_{j>k}\widehat p_{t,k}(j).
\end{align}
Первая величина показывает вероятность того, что лучший момент уже был пропущен. Вторая --- вероятность того, что текущий шаг является оптимальным. Третья --- вероятность того, что полезное уточнение еще впереди. Правило остановки строится поверх этих величин.

